% FTRL-DOR rankings for the LEA case study (adjacent-only likelihood,
% e_max = n_max = 30).  Display scaling:
% U_disp(a) = U(a) / U_max * 100, U_max = sum of all piece weights
% (utility of the ideal alternative scoring 100 on every criterion).

\begin{table}[h!]
\centering
\caption{Ranking of the alternatives considering the parameters estimated through FTRL-DOR after each of the $T=3$ sessions, in the case of linear marginal utility functions. The utilities are expressed as percentages of the utility of the ideal alternative, which scores 100 on every criterion. The reference alternatives are written in bold.}
\label{tab:RANK_lin}
\begin{tabularx}{\textwidth}{c|Cc|Cc|Cc}
\toprule
\multirow{2}{*}{Rank}&\multicolumn{2}{c}{$T=1$}&\multicolumn{2}{c}{$T=2$}&\multicolumn{2}{c}{$T=3$}\\
\cmidrule(lr){2-3}\cmidrule(lr){4-5}\cmidrule(lr){6-7}
&Region&$U\left(a\right)$&Region&$U\left(a\right)$&Region&$U\left(a\right)$\\
\midrule
1 & Veneto & 96.09 & Veneto & 96.46 & Veneto & 96.94 \\
2 & Tuscany & 95.31 & Tuscany & 95.22 & \textbf{Tuscany} & \textbf{95.12} \\
3 & Emilia Romagna & 92.80 & \textbf{Emilia Romagna} & \textbf{93.22} & \textbf{Emilia Romagna} & \textbf{94.10} \\
4 & A. P. of Trento & 92.74 & A. P. of Trento & 92.72 & A. P. of Trento & 93.58 \\
5 & Piedmont & 90.13 & Piedmont & 90.69 & Piedmont & 91.41 \\
6 & \textbf{Lombardy} & \textbf{85.92} & Umbria & 86.66 & \textbf{Lombardy} & \textbf{88.47} \\
7 & Umbria & 85.90 & \textbf{Lombardy} & \textbf{86.64} & Umbria & 88.19 \\
8 & \textbf{Marche} & \textbf{82.29} & \textbf{Marche} & \textbf{80.72} & \textbf{Friuli Venezia Giulia} & \textbf{80.03} \\
9 & Friuli Venezia Giulia & 78.50 & \textbf{Friuli Venezia Giulia} & \textbf{79.27} & \textbf{Marche} & \textbf{78.64} \\
10 & Apulia & 75.79 & Apulia & 74.69 & Apulia & 73.90 \\
11 & Liguria & 72.36 & Liguria & 70.09 & \textbf{Lazio} & \textbf{67.10} \\
12 & Lazio & 71.53 & Lazio & 69.45 & \textbf{Campania} & \textbf{66.08} \\
13 & \textbf{Campania} & \textbf{68.43} & \textbf{Campania} & \textbf{67.54} & Liguria & 66.03 \\
14 & A. P. of Bolzano & 67.17 & A. P. of Bolzano & 66.97 & A. P. of Bolzano & 65.36 \\
15 & Molise & 64.20 & Sardinia & 64.59 & Sardinia & 64.97 \\
16 & Sardinia & 64.09 & Molise & 63.93 & \textbf{Basilicata} & \textbf{63.54} \\
17 & \textbf{Basilicata} & \textbf{63.04} & \textbf{Basilicata} & \textbf{62.81} & Molise & 62.78 \\
18 & Abruzzo & 60.11 & Calabria & 58.76 & Aosta Valley & 62.06 \\
19 & Calabria & 59.08 & Aosta Valley & 57.63 & Calabria & 60.10 \\
20 & Sicily & 57.05 & \textbf{Abruzzo} & \textbf{57.23} & \textbf{Abruzzo} & \textbf{54.93} \\
21 & Aosta Valley & 55.64 & Sicily & 54.02 & \textbf{Sicily} & \textbf{51.32} \\
\bottomrule
\end{tabularx}
\end{table}

\begin{table}[h!]
\centering
\caption{Ranking of the alternatives considering the parameters estimated through FTRL-DOR after each of the $T=3$ sessions, in the case of piecewise linear marginal utility functions ($\gamma_j=3$). The utilities are expressed as percentages of the utility of the ideal alternative, which scores 100 on every criterion. The reference alternatives are written in bold.}
\label{tab:RANK_pw}
\begin{tabularx}{\textwidth}{c|Cc|Cc|Cc}
\toprule
\multirow{2}{*}{Rank}&\multicolumn{2}{c}{$T=1$}&\multicolumn{2}{c}{$T=2$}&\multicolumn{2}{c}{$T=3$}\\
\cmidrule(lr){2-3}\cmidrule(lr){4-5}\cmidrule(lr){6-7}
&Region&$U\left(a\right)$&Region&$U\left(a\right)$&Region&$U\left(a\right)$\\
\midrule
1 & Veneto & 94.48 & Veneto & 93.59 & Veneto & 93.98 \\
2 & Tuscany & 92.57 & Tuscany & 91.05 & \textbf{Tuscany} & \textbf{91.90} \\
3 & A. P. of Trento & 91.92 & A. P. of Trento & 88.30 & A. P. of Trento & 90.81 \\
4 & Emilia Romagna & 90.67 & \textbf{Emilia Romagna} & \textbf{88.28} & \textbf{Emilia Romagna} & \textbf{89.68} \\
5 & Piedmont & 85.42 & Piedmont & 82.89 & Piedmont & 84.12 \\
6 & \textbf{Lombardy} & \textbf{82.46} & \textbf{Lombardy} & \textbf{77.34} & \textbf{Lombardy} & \textbf{80.49} \\
7 & Umbria & 81.11 & Umbria & 76.61 & Umbria & 79.18 \\
8 & \textbf{Marche} & \textbf{69.28} & \textbf{Marche} & \textbf{62.91} & \textbf{Marche} & \textbf{66.60} \\
9 & Friuli Venezia Giulia & 66.54 & \textbf{Friuli Venezia Giulia} & \textbf{61.09} & \textbf{Friuli Venezia Giulia} & \textbf{63.66} \\
10 & Apulia & 62.44 & Apulia & 53.38 & Apulia & 58.84 \\
11 & Liguria & 56.04 & Liguria & 49.72 & Liguria & 52.07 \\
12 & Lazio & 55.20 & Lazio & 45.12 & \textbf{Lazio} & \textbf{50.93} \\
13 & \textbf{Campania} & \textbf{50.42} & A. P. of Bolzano & 42.85 & \textbf{Campania} & \textbf{45.94} \\
14 & A. P. of Bolzano & 49.12 & \textbf{Campania} & \textbf{41.51} & A. P. of Bolzano & 45.29 \\
15 & Molise & 46.43 & Molise & 38.33 & Molise & 42.11 \\
16 & Sardinia & 45.78 & Sardinia & 36.49 & Sardinia & 41.70 \\
17 & \textbf{Basilicata} & \textbf{44.86} & \textbf{Basilicata} & \textbf{34.37} & \textbf{Basilicata} & \textbf{37.42} \\
18 & Abruzzo & 44.81 & Aosta Valley & 33.79 & \textbf{Abruzzo} & \textbf{35.62} \\
19 & Aosta Valley & 42.78 & \textbf{Abruzzo} & \textbf{33.69} & Aosta Valley & 33.75 \\
20 & Sicily & 42.02 & Sicily & 31.04 & \textbf{Sicily} & \textbf{32.29} \\
21 & Calabria & 40.94 & Calabria & 29.91 & Calabria & 31.10 \\
\bottomrule
\end{tabularx}
\end{table}
